\section{Graph-Based Campus Navigation}
\label{sec:nav}
\begin{figure}[!t]
\centering
\begin{tikzpicture}[
  font=\scriptsize, node distance=4mm,
  b/.style={draw, rounded corners=2pt, align=center, text width=6.6cm, minimum height=6mm, inner sep=3pt, fill=black!3},
  pf/.style={-{Latex[length=1.8mm]}}
]
\node[b] (req) {\code{GET /api/v1/navigate} with \code{start\_node\_id}, \code{destination\_node\_id}, \code{accessible\_only}};
\node[b, below=of req] (build) {\code{build\_graph(db, accessible\_only)}: query all \textsf{Node}, all \textsf{Edge}; expand bidirectional edges; drop non-accessible if requested};
\node[b, below=of build] (astar) {\code{find\_shortest\_path}: A$^{\star}$ with binary heap, $f(n){=}g(n){+}h(n)$};
\node[b, below=of astar] (heur) {heuristic $h$: planar $\|\cdot\|_2$ $\rightarrow$ equirectangular great-circle $\rightarrow$ $0$};
\node[b, below=of heur] (fmt) {\code{format\_directions}: template turn-by-turn from \code{edge\_type}/\code{direction}/\code{floor\_transition}};
\node[b, below=of fmt] (resp) {\code{RouteResponse}: node ids, names, distance, walk time, accessibility, steps};
\draw[pf] (req)--(build); \draw[pf] (build)--(astar); \draw[pf] (astar)-- (heur);
\draw[pf] (heur.east) to[bend left=40] (astar.east);
\draw[pf] (astar)--(fmt); \draw[pf] (fmt)--(resp);
\end{tikzpicture}
\caption{A$^{\star}$ navigation pipeline. The graph is rebuilt from the live
relational schema on every request; there is no persistent in-process graph
cache and (currently) no front-end consumer of the endpoint.}
\label{fig:astar}
\end{figure}

\begin{table}[!t]
\centering
\caption{Campus Navigation Approaches and the Implemented Choice.}
\label{tab:nav}
\footnotesize
\renewcommand{\arraystretch}{1.2}
\begin{tabular}{@{}p{1.7cm}p{2.0cm}p{1.5cm}p{1.9cm}@{}}
\toprule
\textbf{Approach} & \textbf{Optimality} & \textbf{Needs coords} & \textbf{Role in \system{}} \\
\midrule
BFS / DFS traversal & unweighted only & no & not used \\
Dijkstra \cite{dijkstra1959note} & optimal (non-neg.\ weights) & no & single-source nearest-facility queries \\
A$^{\star}$ \cite{hart1968astar} & optimal if $h$ admissible & optional & primary point-to-point pathfinder \\
GPS turn-by-turn & n/a & yes (outdoor) & built, then removed from codebase \\
Panorama link traversal & n/a & no & the tour (Section~\ref{sec:tour}), not routing \\
\midrule
\multicolumn{4}{@{}p{7.6cm}@{}}{\textbf{Implemented:} A$^{\star}$ with a three-tier
degrading heuristic; at the zero tier it is provably identical to Dijkstra,
so a single code path covers both. Back end complete; no UI consumer.} \\
\bottomrule
\end{tabular}
\end{table}


The wayfinding engine (\code{backend/app/navigation/}) computes shortest
walking routes over the campus graph; Table~\ref{tab:nav} situates the
implemented A$^{\star}$ choice among alternative routing approaches. It has
been intact and unit-tested since an early project phase; a later scope
revision removed its natural-language chat entry point but kept the engine,
which is now exposed only through \code{GET /api/v1/navigate} and internal
nearest-panorama lookups. It currently has \textbf{no front-end consumer} --- a
previously built map UI with GPS turn-by-turn guidance was deliberately
removed from the codebase.

\subsection{Database-Backed Graph Construction}
The navigation graph is \emph{not} stored as a file; it is rebuilt on every
request from the live relational schema (Algorithm~\ref{alg:graphbuild}).
\code{build\_graph} loads all $\textsf{Node}$ rows and all $\textsf{Edge}$
rows (filtering out non-accessible edges when \code{accessible\_only} is
set), then constructs an adjacency list $A: \mathrm{id} \to
[\textsf{GraphEdge}]$ in which each directed edge $(u,v)$ contributes
$A[u] \mathrel{+}= (v, \text{dist}, e)$ and, if
\code{is\_bidirectional}, also $A[v] \mathrel{+}= (u, \text{dist}, e)$.

\begin{algorithm}[!t]
\caption{Database-Backed Graph Construction}
\label{alg:graphbuild}
\begin{algorithmic}[1]
\Require DB session; flag \code{accessible\_only}
\Ensure \code{CampusGraph}$(A,\ \text{nodesById})$
\State $\text{nodesById} \gets \{\, n.\text{id} \mapsto n : n \in \textsc{SelectAll}(\textsf{Node}) \,\}$
\State $Q \gets \textsc{Select}(\textsf{Edge})$;\ \textbf{if} \code{accessible\_only} \textbf{then} $Q \gets Q\ \textbf{where}\ e.\text{accessible}$
\State $A \gets \{\, i \mapsto [\,] : i \in \text{nodesById} \,\}$
\ForAll{$e \in Q$}
  \State $A[e.\text{src}] \mathrel{+}= (e.\text{tgt},\ e.\text{distance},\ e)$
  \If{$e.\text{is\_bidirectional}$}\ $A[e.\text{tgt}] \mathrel{+}= (e.\text{src},\ e.\text{distance},\ e)$ \EndIf
\EndFor
\State \textbf{return} \code{CampusGraph}$(A,\ \text{nodesById})$
\end{algorithmic}
\end{algorithm}

Using the normalized relational schema as the live source of truth has
three benefits for this project: the graph is always consistent with the
tour and the minimap (which read the same tables); edits go through the
same CRUD/migration tooling as every other campus fact; and accessibility
filtering is a query predicate rather than a separate data set. The cost is
that each route request pays an $O(|V| + |E|)$ rebuild ($180 + 325$ rows);
Section~\ref{sec:eval} proposes measuring whether this warrants a cache.

\subsection{A$^{\star}$ Shortest Path}
The search (Algorithm~\ref{alg:astar}, Fig.~\ref{fig:astar}) is textbook
A$^{\star}$~\cite{hart1968astar,russell2020aima} with a binary heap. It
expands the node minimizing
\begin{equation}
f(n) \;=\; g(n) + h(n, t),
\label{eq:fn}
\end{equation}
where $t$ is the goal node. The accumulated cost uses the stored edge
distance,
\begin{equation}
g(v) \;=\; g(u) + \mathrm{dist}(u,v), \qquad (u,v) \in E,
\label{eq:gn}
\end{equation}
and $\mathrm{dist}(u,v)$ is the \code{Edge.distance} column (a surveyed or
authored float, not derived from coordinates). Path reconstruction walks a
\code{came\_from} map from goal to start. The equal-node fast path returns
immediately; a start or goal absent from the graph raises a typed
no-path error.

\begin{algorithm}[!t]
\caption{A$^{\star}$ Campus Pathfinding}
\label{alg:astar}
\begin{algorithmic}[1]
\Require graph $(A, \text{nodesById})$; start $s$; goal $t$
\Ensure path node ids, edge list, total distance
\If{$s \notin \text{nodesById}$ \textbf{or} $t \notin \text{nodesById}$}\ \textbf{raise} NoPath \EndIf
\If{$s = t$}\ \textbf{return} $([s],\ [\,],\ 0)$ \EndIf
\State open $\gets$ min-heap $\{(0, s)\}$;\ $g[s] \gets 0$;\ visited $\gets \emptyset$;\ came $\gets \{\}$
\While{open $\ne \emptyset$}
  \State $(\cdot,u) \gets \textsc{HeapPop}(\text{open})$
  \If{$u \in$ visited}\ \textbf{continue} \EndIf
  \State visited $\gets$ visited $\cup \{u\}$
  \If{$u = t$}\ \textbf{return} \textsc{Reconstruct}(came, $s$, $t$, $g[t]$) \EndIf
  \ForAll{$(v, w, e) \in A[u]$}
    \State $g' \gets g[u] + w$
    \If{$g' < g.\text{get}(v, \infty)$}
       \State $g[v] \gets g'$;\ came$[v] \gets (u, e)$
       \State \textsc{HeapPush}$\big(\text{open},\ (g' + h(\text{nodesById}[v], \text{nodesById}[t]),\ v)\big)$
    \EndIf
  \EndFor
\EndWhile
\State \textbf{raise} NoPath
\end{algorithmic}
\end{algorithm}

\subsection{Three-Tier Degrading Heuristic}
The heuristic adapts to whatever spatial information the two nodes share:
\begin{equation}
h(a,b) =
\begin{cases}
\sqrt{(a_x - b_x)^2 + (a_y - b_y)^2}, & \text{planar coords on both,}\\[1.2ex]
\sqrt{(\Delta\!X)^2 + (\Delta\!Y)^2}, & \text{geographic coords on both,}\\[1.2ex]
0, & \text{otherwise,}
\end{cases}
\label{eq:heur}
\end{equation}
with the equirectangular approximation
\begin{align}
\Delta\!X &= (a_{\text{lon}} - b_{\text{lon}})\cdot 111{,}320 \cdot \cos\!\Big(\tfrac{a_{\text{lat}} + b_{\text{lat}}}{2}\Big),\\
\Delta\!Y &= (a_{\text{lat}} - b_{\text{lat}})\cdot 110{,}540,
\end{align}
adequate at campus scale. All three tiers are \emph{admissible}: tiers~1
and~2 are straight-line lower bounds on any walk, and $h \equiv 0$ trivially
never overestimates. At tier~3 the priority in~\eqref{eq:fn} reduces to
$f(n) = g(n)$, i.e.\ the search \emph{is} Dijkstra's
algorithm~\cite{dijkstra1959note}. A separate \code{single\_source\_distances}
routine runs plain Dijkstra for nearest-facility queries (nearest room,
nearest panorama), where there is no single fixed goal for a heuristic to
aim at. Because all building \code{latitude}/\code{longitude} values and
most node coordinates are currently null, the system in its present state
runs predominantly at tier~3 (Dijkstra behavior); this is correct, just
uninformed, and improves automatically as real coordinates are surveyed.

\subsection{Turn-by-Turn Directions}
\code{format\_directions} converts the node/edge path into template-based
steps with no LLM involvement: a movement verb keyed on \code{edge\_type}
(``Walk through the corridor'', ``Take the stairs'', ``Take the
elevator''), a turn cue keyed on the relative \code{direction} enumeration
(``Turn left.''), a floor-name cue on floor-transition edges, and, on the
final step into a room, a ``\emph{Room X} is on your left/right'' cue. The
\code{RouteResponse} carries the ordered node ids and names, total
distance, an estimated walk time (summed per-edge \code{walking\_time}, or
distance $/\,1.2\,\text{m\,s}^{-1}$), and a whole-path accessibility flag
(true iff every edge on the path is wheelchair-passable, independent of
whether accessible-only routing was requested).
